\section{Dataset and Evaluation Protocol}
\label{sec:data}

\paragraph{Study cohort.}
We use xView3-SAR Sentinel-1 ground-range-detected scenes with VH and VV
polarizations at 10\,m pixel spacing~\cite{paolo2022xview3}. Scene
identifiers define every split. A frozen seed-0 cohort contains 150 study
scenes: 111 train, 23 development, and 16 test, stratified by coarse
geographic region and shoreline presence. The protocol reserves a separate
set of 50 human-verified scenes for one sealed final evaluation. The four
training budgets are exact nested sets of 12, 28, 56, and 111 scenes:
\begin{equation}
S_{10} \subset S_{25} \subset S_{50} \subset S_{100}.
\end{equation}
Every arm sees the same scene set at a given budget, so a budget never
changes which scenes one initialization sees relative to another. The
preprocessing path converts each scene into a fixed three-channel
decibel-domain tensor $[\mathrm{VH}, \mathrm{VV}, \mathrm{VH}-\mathrm{VV}]$.
It computes normalization statistics from training scenes only. It then cuts
overlapping 800-pixel chips, and optimization samples shared 512-pixel
crops from those chips.

\paragraph{Ground-truth contract.}
One target-semantics contract serves training and evaluation
(Fig.~\ref{fig:gt-contract}, Appendix~\ref{app:figures}). A row is a
positive vessel point only when \texttt{is\_vessel} is true at HIGH or
MEDIUM confidence. HIGH/MEDIUM non-vessels are background. LOW-confidence
rows define ignore regions. Table~\ref{tab:data-protocol} lists the
annotation support of every partition under this contract; a committed
audit receipt binds the counts to the frozen split file by SHA-256. These
support counts are pre-registered integrity checks: evaluation stops if
regenerated counts differ. Training never stages test or human-verified
labels. Appendix~\ref{app:figures} (Fig.~\ref{fig:study-data}) shows the
scale of one study scene against a single chip and the scene centroids of
the frozen splits.

\begin{table}[t]
\centering
\caption{Frozen scene-level data protocol and annotation support under the
corrected ground-truth contract. The 50 verified scenes remain sealed until
the held-out test matrix and review gates complete.}
\label{tab:data-protocol}
\small
\setlength{\tabcolsep}{6pt}
\begin{tabular}{@{}lcccc@{}}
\toprule
Partition & Scenes & Vessel H/M & Non-vessel H/M & LOW ignore \\
\midrule
Training-time dev subset & \HevCountDevEightScenes & \HevCountDevEightPositive & \HevCountDevEightBackground & \HevCountDevEightIgnore \\
Development (full) & \HevCountDevFullScenes & \HevCountDevFullPositive & \HevCountDevFullBackground & \HevCountDevFullIgnore \\
Test & \HevCountTestScenes & \HevCountTestPositive & \HevCountTestBackground & \HevCountTestIgnore \\
Verified final & 50 & \multicolumn{3}{c}{sealed; human verified} \\
\bottomrule
\end{tabular}
\end{table}

\paragraph{Selection and scoring.}
The frozen scorer greedily matches predictions to ground truth within
200\,m and aggregates precision, recall, and F1 across scenes. Every five
epochs, each run performs full-scene inference on the same eight
development scenes and sweeps the score threshold that maximizes
development F1. The run records the best value with the retained
checkpoint. The protocol binds that complete operating point (epoch,
threshold, and counts) to the SHA-256 of the exact checkpoint. Held-out
stages apply each stored threshold verbatim, first to the 16 test scenes
and then, once, to the human-verified set. Held-out annotations never
retune a threshold or select a checkpoint. The frozen split is leakage-free
but not geographically independent: Sentinel-1 revisits place held-out
scenes near training scenes, so we describe absolute performance as
in-region generalization.
